%!TEX root = ../main/main.tex

Responda las siguientes preguntas.
\begin{enumerate}
\item {[3 puntos]} Sea $\mathcal{M}=\mathcal{K}=\mathcal{C}=\{0, 1\}^n$ y considere el esquema criptográfico $(\textit{Gen}, \textit{Enc}, \textit{Dec})$ donde $\textit{Enc}$ y $\textit{Dec}$ son definidos tal como en OTP, pero $\textit{Gen}$ \textbf{no} es la distribución uniforme. Demuestre que este esquema criptográfico no es perfectamente secreto.

\item {[3 puntos]} Formalice la siguiente afirmación y demuestre que un sistema criptográfico la satisface si y sólo si dicho esquema cumple que es perfectamente secreto:
\begin{center}
  \it{La probabilidad de ver un texto cifrado $c\in \mathcal{C}$ sin conocimiento previo es igual a la probabilidad de ver $c$ conociendo el mensaje enviado de antemano.}
  \end{center}
\end{enumerate}

\medskip

\paragraph{Corrección.} La asignación de puntajes para cada pregunta de este problema es la siguiente.
\begin{enumerate}
  \item
    \begin{itemize}
      \item {[1 punto]} Usa una definición más débil de \emph{perfectamente secreto} en la demostración; para dicha definición la demostración es correcta.
      \item {[2 puntos]} Usa la definición correcta y la demostración tiene sentido pero presenta algún problema. Por ejemplo, no cuantifica correctamente los elementos de la demostración, como los textos cifrados o la distribución de probabilidad $\mathbb{D}$ sobre los textos planos.
      \item {[3 puntos]} Usa la definición correcta de \emph{perfectamente secreto} y la demostración es correcta.
    \end{itemize}
  \item
    \begin{itemize}
      \item {[1 punto]} La formalización es correcta, pero la demostración presenta muchos problemas o es evidentemente incorrecta.
      \item {[2 puntos]} La formalización es correcta, pero la demostración es válida en una sola dirección.
      \item {[3 puntes]} La formalización es correcta, y la demostración es válida en ambas direcciones.
    \end{itemize}
\end{enumerate}

\medskip
